\ticket{14}{Байесовские сети как способ представления вероятностного пространства. Построение байесовских сетей, примеры их использования}

\begin{examplan}
    \item Дать определение байесовской сети.
    \item Объяснить разложение совместного распределения и условную независимость.
    \item Описать построение структуры и параметров сети.
    \item Привести примеры использования.
\end{examplan}

\subsection*{Короткая версия для письменного ответа}

\textbf{Байесовская сеть} --- ориентированный ациклический граф, вершины которого соответствуют случайным переменным, а дуги отражают прямые вероятностные зависимости. Каждому узлу задается распределение $P(X_i\mid Parents(X_i))$.

Совместное распределение раскладывается как:
\[
P(X_1,\ldots,X_n)=\prod_{i=1}^{n}P(X_i\mid Parents(X_i)).
\]

\subsection*{Смысл модели}

Байесовская сеть компактно представляет вероятностное пространство за счет условных независимостей. Если переменная имеет родителей, то при известных родителях она считается независимой от непотомков.

Количественная часть сети --- таблицы условных вероятностей для дискретных переменных или параметры условных распределений для непрерывных.

\subsection*{Вероятностный вывод}

Типичные запросы:
\begin{itemize}
    \item вычислить апостериорную вероятность $P(Query\mid Evidence)$;
    \item найти наиболее вероятное объяснение наблюдений;
    \item предсказать ненаблюдаемые переменные по свидетельствам.
\end{itemize}

Алгоритмы вывода бывают точные, например исключение переменных и распространение доверия в деревьях, и приближенные, например MCMC и другие методы Монте-Карло.

\subsection*{Построение сети}

Построение включает два аспекта.

\textbf{Структура сети}:
\begin{enumerate}
    \item выбрать случайные переменные и их значения;
    \item определить прямые зависимости между переменными;
    \item построить ациклический ориентированный граф.
\end{enumerate}

Структура может задаваться экспертно, обучаться из данных или строиться комбинированно.

\textbf{Параметры сети}:
\begin{itemize}
    \item экспертные оценки вероятностей;
    \item оценка по полным данным через частоты и максимальное правдоподобие;
    \item обучение по неполным данным, например EM-алгоритмом;
    \item сглаживание при малом числе наблюдений.
\end{itemize}

\subsection*{Примеры}

\textbf{Медицинская диагностика.} Переменные: сезон, аллергия, простуда, насморк, кашель, температура. По наблюдениям \textit{насморк = да} и \textit{кашель = да} можно вычислить вероятность простуды.

\textbf{Фильтрация спама.} Наивный байесовский классификатор является частным случаем байесовской сети, где класс письма влияет на признаки-слова, а признаки условно независимы при известном классе.

Другие применения: диагностика неисправностей, оценка рисков, биоинформатика, поддержка принятия решений.

\begin{important}
    \item Байесовская сеть состоит из DAG и таблиц условных вероятностей.
    \item Граф кодирует условные независимости.
    \item Совместное распределение раскладывается в произведение локальных условных распределений.
    \item Структура и параметры могут задаваться экспертно или обучаться из данных.
    \item Наивный Байес --- простой пример байесовской сети.
\end{important}
